\section{Simple bias configuration}
Mô phỏng mạch Hình~\ref{fig:3-07} ($V_{CC} = 12\,\text{V}$, $R_1 = 80\,\text{k}\Omega$, $R_2 = 40\,\text{k}\Omega$, $R_E = 3.3\,\text{k}\Omega$, $\beta = 100$) với $R_C = 10\,\Omega$ và $R_C = 1\,\text{k}\Omega$; nhận xét sự thay đổi của $I_E$ và giải thích.

\given{fig3_07}{0.5\linewidth}{Simple bias configuration\label{fig:3-07}}

\subsection{Kết quả mô phỏng}
\simfigs{ex05_rc10}{Bước 1: $R_C = 10\,\Omega$}
\simfigs{ex05_rc1k}{Bước 2: $R_C = 1\,\text{k}\Omega$}

\begin{table}[H]
\centering
\begin{tabular}{l|c|c}
\toprule
 & $R_C = 10\,\Omega$ & $R_C = 1\,\text{k}\Omega$ \\
\midrule
$I_B$ & $8.969\,\mu\text{A}$ & $8.969\,\mu\text{A}$ \\
$I_C$ & $896.9\,\mu\text{A}$ & $896.9\,\mu\text{A}$ \\
$I_E$ & $905.9\,\mu\text{A}$ & $905.9\,\mu\text{A}$ \\
$V_B$ & $3.761\,\text{V}$ & $3.761\,\text{V}$ \\
$V_E$ & $2.989\,\text{V}$ & $2.989\,\text{V}$ \\
$V_C$ & $11.99\,\text{V}$ & $11.10\,\text{V}$ \\
$V_{CE} = V_C - V_E$ & $9.00\,\text{V}$ & $8.11\,\text{V}$ \\
\bottomrule
\end{tabular}
\caption{Kết quả mô phỏng bài 3.5}
\end{table}

\textbf{Nhận xét:} khi $R_C$ tăng từ $10\,\Omega$ lên $1\,\text{k}\Omega$ (gấp 100 lần), $I_E$ \textbf{gần như không đổi} ($905.9\,\mu\text{A}$ ở cả hai trường hợp), chỉ có $V_C$ giảm ($11.99\,\text{V} \to 11.10\,\text{V}$) do sụt áp trên $R_C$ lớn hơn.

\subsection{Phân tích mạch}
Dùng định lý Thevenin ở phía base:
\[ V_{TH} = V_{CC}\frac{R_2}{R_1 + R_2} = 12\cdot\frac{40}{80+40} = 4\,\text{V}, \qquad R_{TH} = R_1\,\|\,R_2 = \frac{80\cdot 40}{80+40} \approx 26.7\,\text{k}\Omega \]
Viết KVL cho vòng base--emitter (với $I_E = (\beta+1)I_B$):
\[ V_{TH} = I_B R_{TH} + V_{BE} + (\beta+1) I_B R_E \;\Rightarrow\; I_B = \frac{V_{TH} - V_{BE}}{R_{TH} + (\beta+1)R_E} \]
\[ I_B = \frac{4 - 0.7}{26.7\,\text{k} + 101\cdot 3.3\,\text{k}} \approx 9.17\,\mu\text{A}, \quad I_C = \beta I_B \approx 0.917\,\text{mA}, \quad I_E = (\beta+1) I_B \approx 0.926\,\text{mA} \]
\[ V_E = I_E R_E \approx 3.06\,\text{V}, \qquad V_B = V_E + V_{BE} \approx 3.76\,\text{V} \]
Biểu thức của $I_B$, $I_C$, $I_E$ \textbf{không chứa $R_C$}, nên đổi $R_C$ không làm $I_E$ thay đổi. Chỉ điện áp collector phụ thuộc $R_C$:
\[ V_C = V_{CC} - I_C R_C \approx \begin{cases} 12 - 0.917\,\text{mA}\cdot 10\,\Omega \approx 11.99\,\text{V} & (R_C = 10\,\Omega)\\ 12 - 0.917\,\text{mA}\cdot 1\,\text{k}\Omega \approx 11.08\,\text{V} & (R_C = 1\,\text{k}\Omega)\end{cases} \]
Cả hai trường hợp $V_{CE} \approx 8.9\,\text{V}$ và $8.0\,\text{V}$ vẫn rất lớn hơn $V_{CE(sat)}$, nên BJT luôn ở vùng tích cực; trong vùng này $I_C = \beta I_B$ do dòng base quyết định chứ không do $R_C$.

\begin{notebox}{Giải thích hiện tượng}
Điện trở $R_E$ tạo hồi tiếp âm, cùng với cầu chia áp $R_1$--$R_2$ giữ cố định điện áp $V_B$: nếu $I_E$ có xu hướng tăng thì $V_E = I_E R_E$ tăng, $V_{BE} = V_B - V_E$ giảm, làm $I_B$ và $I_E$ giảm trở lại. Nhờ đó điểm làm việc ổn định, $I_E$ chỉ phụ thuộc $V_{TH}$, $R_{TH}$, $R_E$, $\beta$, $V_{BE}$ mà không phụ thuộc $R_C$ (miễn là BJT còn ở vùng tích cực). $I_E$ mô phỏng thấp hơn lý thuyết khoảng $2\%$ do $V_{BE}$ thực là $0.772\,\text{V}$.
\end{notebox}
